\Needspace{35\baselineskip}\SourceLeaf{84}

\begin{center}
I.
\end{center}

\noindent\hspace*{\parindent}Da für $u = 0$, $x = t$ wird, so ist nach \textsc{Taylor}s Lehrsatze

\[
\varphi x = \varphi t + u \frac{d\varphi x}{du} + \frac{uu}{1.2}\frac{d^2\varphi x}{du^2}
+ \frac{u^3}{1.2.3}\frac{d^3\varphi x}{du^3} + \text{etc.},
\]

\Needspace{23\baselineskip}\noindent wo der Werth der Grössen $\dfrac{d\varphi x}{du}$, $\dfrac{d^2\varphi x}{du^2}$ etc.\ für
$u = 0$ zu nehmen ist. Allein es ist allgemein

\[
\frac{d\varphi x}{du} = X\frac{d\varphi x}{dt},
\]

\[
\frac{d^2\varphi x}{du^2} = \frac{dX^2 \frac{d\varphi x}{dt}}{dt},
\]

\[
\frac{d^3\varphi x}{du^3} = \frac{d^2X^3 \frac{d\varphi x}{dt}}{dt^2},
\]

\[
\cdot\quad\cdot\quad\cdot\quad\cdot\quad\cdot\quad\cdot\quad\cdot\quad\cdot\quad\cdot,
\]

\noindent wie ich sogleich beweisen werde. Ferner wird man leicht einsehen, dass für $u = 0$, d.\,i.\
$x = t$ diese Ausdrücke

\[
X\frac{d\varphi x}{dt},\quad \frac{dX^2 \frac{d\varphi x}{dt}}{dt},\quad
\text{etc.\ in}\quad X\frac{d\varphi x}{dx},\quad \frac{dX^2 \frac{d\varphi x}{dx}}{dx},\quad \text{etc.}
\]

\noindent übergehen, worin dann $x = t$ zu setzen ist. Dadurch wird also

\[
\varphi x = \varphi t + uX\frac{d\varphi x}{dx} + \frac{uu}{1.2}\frac{dX^2 \frac{d\varphi x}{dx}}{dx}
+ \frac{u^3}{1.2.3}\frac{d^2X^3 \frac{d\varphi x}{dx}}{dx^2} + \cdot\,\cdot\,\cdot,
\]

\noindent wo $x = t$ zu setzen ist, und dieses ist die Formel des Herrn \textsc{de la Grange}.

\medskip

\Needspace{10\baselineskip}\begin{center}
II.
\end{center}

\noindent\hspace*{\parindent}Die hierbei gebrauchten Verwandlungen beweise ich so:

\noindent\hspace*{\parindent}Ich differentiire die vorgegebene Gleichung $x = t + uX$ nach $t$ und nach $u$,
woraus ich erhalte

\[
\frac{dx}{dt} = 1 + u\frac{dX}{dt} = 1 + u\frac{dX}{dx}\frac{dx}{dt},
\]

\[
\frac{dx}{du} = X + u\frac{dX}{du} = X + u\frac{dX}{dx}\frac{dx}{du}.
\]
